\chapter{Конструкторский раздел}
\label{cha:design}

В данном разделе пишем про алгоритмы, архитектуру ПО и т.п.

Формула \eqref{F:F1} совершено бессмысленна.

\begin{equation}
a= cb
\label{F:F1}
\end{equation}

А формула~\eqref{eq:fourierrow} имеет некоторый смысл.

\begin{eqndesc}
    \begin{equation}\label{eq:fourierrow}
        f(x) = \frac{a_0}{2} + \sum\limits_{k=1}^{+\infty} A_k\cos\left(k\frac{2\pi}{\tau}x+\theta_k\right)
    \end{equation}

    где $A_k$ "--- амплитуда  k-го гармонического колебания,\\
    $A_k$ "--- амплитуда $k$-го гармонического колебания,\\
    $ k\frac{2\pi}{\tau} = k\omega$ "--- круговая частота гармонического колебания,\\
    $\theta_k$ "--- начальная фаза $k$-го колебания.
\end{eqndesc}


Окружение \texttt{cases} опять работает (см. \eqref{F:F2}), спасибо И. Короткову за исправления..


\begin{equation}
a= \begin{cases}
 3x + 5y + z, \mbox{если хорошо} \\
 7x - 2y + 4z, \mbox{если плохо}\\
 -6x + 3y + 2z, \mbox{если совсем плохо}
\end{cases}
\label{F:F2}
\end{equation}

\section{Подсистема всякой ерунды}

Культурная вставка dot-файлов через утилиту dot2tex (рис.~\ref{fig:fig02}).

\begin{figure}
  \centering
% [width=0.5\textwidth] --- регулировка ширины картинки
%   \includegraphics[width=.5\textwidth]{inc/dot/cow2}
  \caption{Рисунок}
  \label{fig:fig02}
\end{figure}


\section*{Выводы}

Что сделали в этом разделе, получили в результате, кратко.
